\chapter{Network Analysis and Insights}
\label{ch:analysis}
Analysis separates five projections: researcher coauthorship, organization
collaboration, directed within-corpus citations, paper--topic membership and
organization--topic contributions. Metrics come from analysis of the graph,
computed offline with NetworkX \citep{networkxDocs}. Each projection includes
its eligible nodes and isolates; collaboration graphs are undirected and
citation arrows run from citing to cited paper.

\section{Publication and topic trends}
The official counts in Figure~\ref{fig:publications} show about 80\% more
papers in 2025 than in 2021, with a decline in 2022. The expansion is uneven
and specific to the selected main-track proceedings.

The topic view asks a different question: how is the labeled part of that
corpus distributed across model-defined research areas? A paper with $L$
accepted labels contributes $1/L$ to each. Annual shares divide the resulting
mass by the number of labeled papers, so a multi-label paper is not counted as
several whole papers. Figure~\ref{fig:topics} combines both venues and displays
the eight topics with the greatest total fractional mass.

\plot{topic-trends}{Inferred topic trends for the eight largest topics by total
fractional mass. A labeled paper contributes total mass one across its labels;
annual denominators combine labeled AAAI and IJCAI papers. Panels share the same
vertical scale.}{fig:topics}
\FloatBarrier

Representation learning is prominent across the period, while the inferred
multimodal share increases in both venues. In the venue-specific frozen tables,
AAAI's multimodal share rises from 6.96\% in 2021 to 16.42\% in 2025;
IJCAI's rises from 7.43\% to 12.90\%. Seeing the same direction in both series
makes this a useful pattern to inspect beyond the changing venue mixture.
However, it remains a shift in classifier-assigned shares, conditional on
labeled papers. Abstentions and the topic-evaluation limits in
Chapter~\ref{ch:quality} constrain a claim about the entire research landscape.

\section{Researcher collaboration network}
The coauthorship projection connects \Researchers{} researchers through
\CoauthorPairs{} distinct pairs. Degree counts different collaborators, while
fractional strength accumulates team-size-adjusted support. Mean degree is
5.00 and the maximum is 90. The low mean alongside a long upper tail indicates
that a relatively small set of records has much broader collaboration reach
than the typical researcher.

\plot{researcher-structure}{Researcher collaboration structure. Left: fraction
of all 32,189 researchers with degree at least $k$; positive degrees are plotted
on log axes, with isolates retained in the denominator. Right: all 2,936
components ranked by size, including isolates; vertical scale is logarithmic.
The curve alone does not establish a power-law model.}{fig:structure}

The largest connected component contains \LargestComponent{} researchers,
or 64.5\% of the projection. There are \ResearcherComponents{} components in
all and \ResearcherIsolates{} isolated researchers. A path through the largest
component means a chain of coauthorships, not that every member has worked
directly with every other member. The contrast between one large component and
many small groups indicates a connected core alongside substantial local
fragmentation.

The time dimension adds context. Table~\ref{tab:annual} counts researchers with
at least one accepted authorship in each year, and assigns each coauthor pair
to the earliest year of its supporting papers. Active researchers can recur in
several rows. The last column therefore measures newly observed pairs within
this five-year window, not all pairs active that year or first collaboration
in a person's career.

% Generated from frozen v4; edit generator.
\begin{table}[!htbp]\centering\small
\caption{\lang{Annual activity in the accepted coauthorship network.}{Hoạt động hằng năm trong mạng đồng tác giả được chấp nhận.}}\label{tab:annual}
\begin{tabularx}{\textwidth}{rYY}
\toprule \lang{Year}{Năm} & \lang{Active researchers}{Nhà nghiên cứu hoạt động} & \lang{First-observed pairs}{Cặp xuất hiện lần đầu}\\\midrule
2021 & \num{6337} & \num{12599}\\
2022 & \num{6136} & \num{10730}\\
2023 & \num{6941} & \num{12998}\\
2024 & \num{9361} & \num{17447}\\
2025 & \num{13082} & \num{26708}\\
\bottomrule\end{tabularx}\end{table}

\FloatBarrier

Active researcher records increase from 6,337 in 2021 to 13,082 in 2025;
first-observed pairs increase from 12,599 to 26,708. Together with the
publication trend, this shows that the later corpus contains both more papers
and a wider accepted collaboration network. Identity withholding can remove
links unevenly, so the graph cannot distinguish every real change in research
activity from metadata coverage effects.

\section{Organization collaboration network}
The organization projection contains \Organizations{} organizations and
\OrganizationPairs{} unique collaboration pairs. It has
\OrganizationIsolates{} isolates, mean degree 12.73 and maximum degree 257.
An edge requires paper-specific affiliations of two distinct authors to two
different organizations. This gives a co-publication relation with a concrete
construction rule.

\textbf{Unique pairs and pair--paper supports differ.} An organization pair is
one graph edge even if it appears on many papers. Its support count records the
number of distinct papers. A paper with several eligible pairs contributes one
support to each pair, and fractional weights additionally divide its total mass
across those pairs. The sector matrix in Figure~\ref{fig:sectors} sums
pair--paper supports, not unique papers or unique graph edges. Its mirrored
cells show the same cross-sector supports and should not be added twice.

\plot{sector-mixing}{Organization collaboration by sector. Cells count distinct
organization-pair by core-paper supports. The symmetric matrix retains unknown
sectors; each eligible pair requires two distinct authors with aligned
paper-specific affiliations.}{fig:sectors}
\FloatBarrier

Academia--academia has the largest displayed cell, with 10,956 supports.
Academia--industry has 3,533. This indicates that academic co-publication is the
largest visible sector pattern while a substantial cross-sector connection is
also present. Sector is an OpenAlex organization property; unknown metadata is
kept explicit. The counts describe publication support rather than contractual
partnerships.

\Needspace{17cm}
Figure~\ref{fig:orgpairs} focuses on the twelve academia--industry pairs with the
most distinct supporting papers across all topics. It turns a dense graph into
a readable comparison of repeated co-publication. This edge ranking contains
eleven unique organizations, so institutions can occur in several rows. It is
not an induced network on those organizations.

\plot{organizations-network}{Twelve leading academia--industry pairs by
accepted distinct-paper support across all topics. Ties use stable endpoint IDs.
Repeated organizations are intentional; rows are unique pairs.}{fig:orgpairs}
\FloatBarrier

A topic-specific view reveals a more focused example. Shanghai Jiao Tong
University--Tencent has 14 distinct supporting papers labeled
\emph{Representation learning}, tied with Xiamen University--Tencent and
Zhejiang University--Alibaba. Two source-checked Shanghai Jiao Tong--Tencent
papers are \href{https://doi.org/10.1609/aaai.v36i1.19954}{\emph{Exploiting
Fine-Grained Face Forgery Clues via Progressive Enhancement Learning}} and
\href{https://doi.org/10.1609/aaai.v36i3.20193}{\emph{DIRL: Domain-Invariant
Representation Learning for Generalizable Semantic Segmentation}} (2022).
Accepted observations align the two organizations to distinct authors. The
14-paper result uses a topic filter, whereas Figure~\ref{fig:orgpairs} uses
all topics.

Organization portfolios also depend on the denominator. For inferred
Representation learning, Zhejiang University has 156 supported papers among
489 organization-associated topic-labeled papers; Peking University has
139 among 381. Their fractional portfolio shares are 25.1\% and 28.8\%.
These shares divide each labeled paper across both its supported organizations
and its labels. They are consequently different from the raw ratios 156/489
and 139/381. Within the observed portfolio, Zhejiang has more supporting papers,
while Peking allocates a greater fractional share to that topic. This is a useful
comparison of volume and specialization, conditional on retained affiliations
and inferred labels.

\section{Community structure and researcher centrality}
Community detection finds relatively dense collaboration groups. The selected
Louvain run uses seed 17 and resolution 0.8 \citep{blondel2008communities}.
Figure~\ref{fig:communities} shows inferred-topic profiles for twelve large
communities. Papers are deduplicated within each community and divide their
mass across accepted topic labels. One paper can be associated with several
communities through its authors.

\plot{community-topic-profiles}{Topic profiles of the first twelve communities
in the stored size order. The ten topics with greatest summed fractional mass
in these profiles are displayed. Each community stores only its top five topics:
grey means unreported, not zero. Community IDs belong to this fixed partition.}{fig:communities}
\FloatBarrier

The profiles show different mixtures rather than a one-community/one-topic
mapping. Representation learning occurs across several groups, while other
inferred topics are more concentrated in particular profiles. Such overlap is
useful for finding neighboring research groups with related interests.
Community boundaries depend on the graph and Louvain settings, so they are
analytical groupings rather than official departments or stable disciplines.

Centrality provides a complementary view. Degree measures the breadth of direct
collaboration; strength reflects repeated, fractionally weighted collaboration.
Betweenness measures how often a researcher lies on shortest paths between
others \citep{freeman1977betweenness}. The frozen analysis uses inverse positive
fractional strength as path distance and approximates weighted betweenness with
256 pivots. Table~\ref{tab:bridges} adds community context to the first five
records instead of interpreting rank alone.

% Generated from frozen v4; edit generator.
\begin{table}[!htbp]\centering\small
\caption{\lang{Leading approximate weighted-betweenness records and community context.}{Các định danh đứng đầu theo độ trung tâm trung gian xấp xỉ có trọng số và cộng đồng tương ứng.}}\label{tab:bridges}
\begin{tabularx}{\textwidth}{Yrrr}
\toprule \lang{Researcher}{Nhà nghiên cứu} & \lang{Betweenness}{\makecell[r]{Độ trung tâm\\trung gian}} & \lang{Community}{Cộng đồng} & \lang{Members}{Số thành viên}\\\midrule
Bo An & 0.042193 & 00126 & \num{367}\\
Lei Feng & 0.034010 & 00126 & \num{367}\\
Hau Chan & 0.021608 & 00000 & \num{687}\\
Shirui Pan & 0.020971 & 00043 & \num{481}\\
Chengjie Wang & 0.019882 & 00058 & \num{953}\\
\bottomrule\end{tabularx}\end{table}


Bo An (A5017743551) ranks first at 0.042193 among 32,189 researchers. This
identity belongs to a 367-member community linked to 206 distinct papers,
including 180 with topic labels. Representation learning accounts for 63 papers
and 52.3333 fractional topic units, or 29.1\% of its fractional topic mass.
Lei Feng, the second-ranked record, belongs to the same selected community.
This illustrates how several bridge records can arise within one broader
collaboration setting rather than representing unrelated isolated leaders.

Two source-checked participations for Bo An are
\href{https://doi.org/10.1609/aaai.v36i8.20851}{\emph{GearNet: Stepwise Dual
Learning for Weakly Supervised Domain Adaptation}} (2022) and
\href{https://doi.org/10.1609/aaai.v37i6.25871}{\emph{Partial-Label Regression}}
(2023). These demonstrate accepted participation; they do not explain the
betweenness score on their own.

\plot{researchers-network}{Bo An (gold) and twelve accepted one-hop neighbors,
selected by descending paper support, fractional weight and stable ID. All
accepted edges among the thirteen selected IDs are shown: 24 edges. Blue edges
touch Bo An and vary in width with support; grey edges connect selected neighbors.
Layout seed 17; this is a bounded neighborhood, not the full ego network.}{fig:ego}
\FloatBarrier

The neighborhood includes links among selected coauthors, not only a star
centered on Bo An. It therefore helps inspect the local collaboration setting,
while global brokerage remains a property of the complete projection.
Interpretation also depends on the support rule: the frozen minimum-two-paper
filter retains only 7.78\% of researcher collaboration edges and has zero
members in common with the base betweenness top twenty. The filter represents
repeated collaboration and strongly fragments the network. This is a material
reason to treat bridge rankings as structural inspection leads. Betweenness
measures paths under the chosen graph definition, not research quality or
real-world influence.

\section{Within-corpus citation network}
The directed citation projection includes all \Papers{} papers and
\CitationEdges{} accepted within-corpus citation edges. A reference to a paper
outside these proceedings is excluded; ambiguous target mappings and self-loops
are dropped, and repeated references to the same target collapse. The graph
has \CitationIsolates{} papers with neither incoming nor outgoing in-corpus
citations, mean total degree 1.07 and maximum in-degree 51.

\plot{top-cited-papers}{Ten highest in-degree papers in the captured citation
network. Each count is the number of distinct citing papers inside this corpus.
Official titles remain unchanged. Counts exclude external citations and do not
normalize publication age.}{fig:cited}
\FloatBarrier

\emph{Informer: Beyond Efficient Transformer for Long Sequence Time-Series
Forecasting} (AAAI 2021, \texttt{paper:aaai:2021:17325}) leads with 51 in-corpus
citing papers. It is an example of a time-series forecasting paper attracting
repeated attention within the selected AI proceedings. The ranking can help a
reader locate frequently referenced methodological starting points, but it is
not a field-wide citation ranking.

The many citation isolates are also meaningful: restricting both endpoints to
two venue series removes most of the wider reference network. An isolated paper
here can still cite, and be cited by, substantial literature elsewhere. Older
papers also have more time to collect citations. Reading the ranking alongside
these two boundaries avoids interpreting a low in-corpus count as low overall
attention.

\section{Key findings}
\begin{enumerate}
\item \textbf{The selected publication landscape expands unevenly.} Output
rises from 2,240 to 4,042 papers after a 2022 dip. Multimodal share rises in both
venues' labeled subsets, suggesting a useful topic trend to inspect alongside
the broader increase in volume.
\item \textbf{Coauthorship combines a large core and many local groups.}
The largest component holds 20,776 of 32,189 researchers, while the graph has
2,936 components. This supports both broad network navigation and local
group-level exploration.
\item \textbf{Cross-sector co-publication is recurrent.} Academia--industry
accounts for 3,533 pair--paper supports. The 14-paper Representation-learning
examples identify concrete repeated publication links for research-program
planning, within the retained affiliation evidence.
\item \textbf{Bridge position and collaboration breadth answer different
questions.} Bo An leads approximate betweenness at 0.042193, but the repeated-support
filter changes the entire top twenty. A shortlist needs its projection and
support rule alongside its names.
\item \textbf{Citation attention is concentrated within a narrow window.}
Informer has the highest in-degree, 51, while 7,302 papers are citation isolates.
Both results reflect the two-venue citation boundary and captured references.
\end{enumerate}
